% Bewertungspaket zum Gesamttest «Betragsfunktionen» — Quelle des PDFs (python3 scripts/build-lp-pdf.py betrag).
% Ganze Punkte je Teilschritt; Werte mit python3 nachgerechnet (05.10.2026, Folgefehler-Fälle eingeschlossen).
\documentclass[11pt]{article}
\usepackage{../lp-druck}
\renewcommand{\lpname}{Betragsfunktionen}
\renewcommand{\lpbereich}{Schwerpunktfach}
\renewcommand{\lpteilgebiet}{3.6}
\renewcommand{\lpfuss}{Bewertungspaket}
\newcolumntype{P}{>{\centering\arraybackslash}p{10mm}}
\newenvironment{raster}{\par\smallskip\noindent\tabularx{\linewidth}{@{}>{\raggedright\arraybackslash}p{38mm}P X@{}}\toprule
  \small\textsc{Teilschritt} & \small\textsc{P} & \small\textsc{Musterlösung}\\\midrule}{\bottomrule\endtabularx\par}
\newcommand{\fehler}[1]{\par\smallskip{\small\textcolor{lprot}{\bfseries Typische Fehler:} #1}}
\newcommand{\E}{\,{\footnotesize\bfseries\color{lpgruen}(E)}}
\newcommand{\A}{\,{\footnotesize\bfseries\color{lpblau}(A)}}
\begin{document}
\lpkopf{Bewertungspaket zum Gesamttest}{}

\begin{lpachtung}{Erst lösen, dann öffnen}
Dieses Paket enthält die Musterlösung. Wer es vor dem Lösen liest, misst am Ende nur, wie gut das Abschreiben gelingt.
\end{lpachtung}

\section*{1 · So gehst du vor}
\begin{enumerate}[leftmargin=*,itemsep=2pt]
\item \textbf{Gesamttest lösen} — vollständig, mit Rechenweg, ohne dieses Paket und ohne Taschenrechner.
\item \textbf{Fotografieren oder scannen}, gut lesbar, eine Seite pro Bild. \textbf{Kein Name} auf den Bildern.
  Handyfotos tragen oft unsichtbar Standort, Zeit und Gerät mit (Metadaten). Lade darum lieber einen Scan oder ein
  Bildschirmfoto des Fotos hoch, oder entferne vorher die Standortangaben.
\item \textbf{Bewerten lassen.} Ob du dafür eine KI nutzt und welche, entscheidest du selbst und in eigener Verantwortung —
  je nachdem, was dir aufgrund deines Alters und deiner persönlichen Situation erlaubt ist (Altersgrenzen und
  Nutzungsbedingungen des Anbieters, allenfalls Einverständnis deiner Eltern). Lade nur deine Lösung und dieses PDF hoch,
  keine persönlichen Daten, und füge den Auftrag aus Abschnitt~2 ein. Ohne KI kannst du dich mit Abschnitt~3 selbst bewerten.
\item \textbf{Rückmeldung prüfen.} Stimmt, was die KI aus deiner Handschrift gelesen hat? Eine KI kann sich irren —
  im Zweifel gilt das Raster in Abschnitt~3.
\item \textbf{Gezielt wiederholen} nach Abschnitt~4.
\end{enumerate}

\needspace{14\baselineskip}
\section*{2 · Auftrag an die KI}
\begin{tcolorbox}[colback=lplinie!25,colframe=lplinie,boxrule=0.5pt,arc=1mm,fontupper=\small]
Du bist Korrektorin oder Korrektor für Mathematik an einer Schweizer Berufsmaturitätsschule. Im Anhang findest du (1)~das Bewertungspaket zum Gesamttest «Betragsfunktionen» mit Musterlösung und Punkteraster und (2)~Fotos meiner handschriftlichen Lösung.\par\medskip
Geh so vor:\par
1. Schreib zu jeder Aufgabe G1 bis G7 zuerst ab, was du in meiner Lösung liest — Rechenweg und Ergebnis. Was du nicht sicher lesen kannst, markierst du mit [unsicher]. Nicht raten.\par
2. Vergib die Punkte genau nach dem Raster im Bewertungspaket (Abschnitt 3), ganze Punkte je Zeile. Ziehe einen Fehler nur einmal ab: Rechne ich mit einem falschen Zwischenergebnis richtig weiter, gibt es die Folgepunkte. Ausnahme: Zeilen mit \textbf{(E)} gibt es nur für das richtige Ergebnis, nie als Folgepunkt. Die Zeilen «Typische Fehler» sagen, welche Rasterzeile bei diesem Fehler 0 Punkte gibt — sie sind kein zusätzlicher Abzug.\par
3. Andere richtige Lösungswege zählen voll. Gleichwertige Schreibweisen zählen voll: Dezimalpunkt oder Dezimalkomma, Bruch oder Dezimalzahl ($\tfrac18 = 0.125$), $y = \ldots$ oder $f(x) = \ldots$, Fallunterscheidung mit $\ge$/$<$ oder $>$/$\le$ an der Grenze, $L = \{-2;\, 4\}$ oder «$x = -2$ oder $x = 4$», Intervall oder Ungleichungskette — ausser wo das Raster eine bestimmte Form verlangt. Drei Stufen: Zeilen mit \textbf{(A)} verlangen nur Ablesen, diese Punkte gibt es auch ohne Rechenweg. Zeilen mit \textbf{(E)} sind Ergebnispunkte: Das Ergebnis muss stimmen \emph{und} aus einem erkennbaren Weg hervorgehen, ausser die Zeile sagt ausdrücklich etwas anderes. Alle übrigen Zeilen bewerten den Weg selbst.\par
4. Begründe jeden Abzug in einem Satz und nenne die Fehlerart (zum Beispiel «Vorzeichen des Knicks vertauscht» oder «nur ein Fall gerechnet»). Schreib die Musterlösung nicht ab — zeig mir, wo mein Fehler liegt.\par
5. Gib zum Schluss eine Tabelle «Aufgabe | Punkte | Begründung», die Gesamtpunktzahl (von 25) und — nach der Selbsteinschätzung in Abschnitt 4 — welches Kapitel des Leitprogramms ich wiederholen soll. Keine Note.\par
6. Fehlt ein Foto oder ist eine Aufgabe unlesbar, sag es, statt sie zu bewerten.
\end{tcolorbox}

\needspace{16\baselineskip}
\section*{3 · Musterlösung und Punkteraster}
\textbf{Allgemein:} Ganze Punkte je Zeile. Ein Fehler wird nur einmal abgezogen (Folgefehler). Gleichwertige Wege und
Schreibweisen zählen voll. In der Spalte P:
\textbf{(A)}~= nur ablesen, zählt auch ohne Rechenweg · \textbf{(E)}~= Ergebnispunkt, Ergebnis richtig
\emph{und} Weg erkennbar (ausser die Zeile sagt ausdrücklich etwas anderes) · ohne Zeichen = die Zeile bewertet den Weg.
Der ganze Test ist ohne Taschenrechner gestellt.
«Typische Fehler» nennt, welche Zeile 0 Punkte gibt, und die Punkte, die bleiben.
\textbf{Teilrichtige Zeilen:} Verlangt eine Zeile mehrere Angaben und stimmt nur ein Teil, gibt die Zeile 0 Punkte —
ausser die Zeile nennt ausdrücklich eine Aufteilung.

\begin{lpaufgabe}{G1}{Die Betragsfunktion}{3 P}
\begin{raster}
(a) & 1\E & $|-7| - 4 = 7 - 4 = 3$\\
(b) & 1 & $|x| = x$ für $x \ge 0$, $-x$ für $x < 0$ (gleichwertig: $x > 0$ / $x \le 0$); Skizze: V mit Knick $(0 \mid 0)$ durch $(\pm 3 \mid 3)$. Fälle und Skizze nötig.\\
(c) & 1\A & $x = 4.5$ oder $x = -4.5$. Beide nötig; zählt auch ohne Rechenweg.\\
\end{raster}
\fehler{(a) $|3 - 10| = 13$ gerechnet: $9$. $-|-4| = 4$ gerechnet: $11$. Beides: $17$. Jeweils Zeile (a) 0.
(b) Fall «$x < 0$: $|x| = x$» oder V unter der $x$-Achse: Zeile (b) 0.
(c) Nur $x = 4.5$: Zeile (c) 0 $\to$ 2 von 3\,P.}
\end{lpaufgabe}

\begin{lpaufgabe}{G2}{Ein Dach skizzieren}{4 P}
\begin{raster}
Knick & 1\A & $(-1 \mid 4)$\\
Steigungen & 1\A & links $+2$, rechts $-2$ (Dach, nach unten offen). Beide nötig.\\
Skizze & 1 & Dach mit Spitze $(-1 \mid 4)$, durch $(0 \mid 2)$ und $(-2 \mid 2)$, Äste als Geraden. Toleranz: ein halbes Kästchen.\\
Nullstellen & 1\E & $-2|x + 1| + 4 = 0 \Rightarrow |x + 1| = 2 \Rightarrow x = 1$ oder $x = -3$\\
\end{raster}
\fehler{Knick bei $(1 \mid 4)$ (Vorzeichen von $u$): Knick-Zeile 0; Skizze als Folgefehler zählt, die Nullstellen ($3$, $-1$) als (E) nicht $\to$ höchstens 2 von 4\,P.
V statt Dach (Minus vor dem Betrag übersehen): Steigungs- und Skizzen-Zeile 0; dann gibt es keine Nullstellen — auch diese Zeile 0.
Nur eine Nullstelle: Nullstellen-Zeile 0.}
\end{lpaufgabe}

\begin{lpaufgabe}{G3}{Gleichung aus dem Graphen}{3 P}
\begin{raster}
$u$, $v$ & 1\A & Knick $(-2 \mid -1)$: $u = -2$, $v = -1$. Beide nötig.\\
$a$ & 1 & mit einem abgelesenen Punkt begründet, z.\,B. $(0 \mid 0)$: vom Knick zwei nach rechts, eine hinauf, $a = \tfrac12 = 0.5$. Ohne Punkt keine Punktzahl.\\
Gleichung & 1\E & $y = 0.5\,|x + 2| - 1$ (gleichwertig: $\tfrac12|x + 2| - 1$)\\
\end{raster}
\fehler{$u = 2$ (Vorzeichen), also $0.5\,|x - 2| - 1$: Zeilen $u$/$v$ und Gleichung 0 $\to$ 1 von 3\,P.
$a = 2$ (Kehrwert der Steigung): Zeilen $a$ und Gleichung 0 $\to$ 1 von 3\,P.}
\end{lpaufgabe}

\begin{lpaufgabe}{G4}{Umklappen}{4 P}
\begin{raster}
(a) Nullstellen & 1\E & $x^2 - 2x - 3 = (x + 1)(x - 3) = 0$: $x = -1$, $x = 3$ (auch mit der Lösungsformel)\\
(a) Scheitel & 1\E & $(1 \mid -4)$ (Mitte der Nullstellen, $f(1) = -4$)\\
(b) Skizze & 1 & Parabel $f$ und darüber $|f|$: ausserhalb von $[-1;\, 3]$ gleich, dazwischen nach oben geklappt — ein W, nie unter der $x$-Achse.\\
(b) Knicke, Scheitel & 1 & Knicke an den Nullstellen aus (a), umgeklappter Scheitel $(1 \mid 4)$. Gilt auch als Folgefehler zu falschen Werten aus (a), wenn die Knicke dort liegen und der Scheitel richtig hochgeklappt ist.\\
\end{raster}
\fehler{Ganze Parabel gespiegelt ($-f$ statt $|f|$): Skizzen-Zeile 0.
Scheitel bei $(1 \mid -4)$ belassen: Die Skizze liegt dann unter der $x$-Achse — Skizzen- und letzte Zeile 0 $\to$ höchstens 2 von 4\,P.
Nullstellen $1$ und $-3$ (Vorzeichen beim Faktorisieren): Zeile (a) Nullstellen 0, Scheitel $(-1 \mid -4)$ ebenfalls 0; (b) als Folgefehler bewerten $\to$ höchstens 2 von 4\,P.}
\end{lpaufgabe}

\begin{lpaufgabe}{G5}{Abschnittsweise schreiben}{4 P}
\begin{raster}
(a) Grenze & 1 & $2x - 5 = 0 \Rightarrow x = 2.5$\\
(a) Terme & 1\E & $|2x - 5| = 2x - 5$ für $x \ge 2.5$, $-2x + 5$ für $x < 2.5$. Beide nötig.\\
(b) Terme & 1\E & $-2x + 3$ für $x < -1$, $5$ für $-1 \le x \le 4$, $2x - 3$ für $x > 4$. Alle drei nötig.\\
(b) Boden & 1\E & Höhe $4 - (-1) = 5$\\
\end{raster}
\fehler{(a) Links $2x + 5$ (nur eine Zahl umgedreht): Terme-Zeile 0.
(a) Grenze bei $5$: Grenz-Zeile 0; die Terme mit falscher Grenze ebenfalls 0.
(b) Nur der Boden, ohne die äusseren Terme: Terme-Zeile (b) 0 $\to$ 3 von 4\,P. Höhe $3$ (aus $|{-1} + 4|$): Boden-Zeile 0.}
\end{lpaufgabe}

\begin{lpaufgabe}{G6}{Gleichung und Ungleichung}{4 P}
\begin{raster}
(a) Fälle & 1 & $3x - 6 = 9$ oder $3x - 6 = -9$ (gleichwertig: $|x - 2| = 3$)\\
(a) Lösung & 1\E & $L = \{-1;\, 5\}$\\
(b) Skizze & 1 & V mit Knick $(-1 \mid 0)$ und Waagrechte $y = 2$; Schnittstellen $-3$ und $1$ markiert. Eine rechnerische Begründung «$x + 1 \le -2$ oder $x + 1 \ge 2$» zählt gleich.\\
(b) Lösung & 1\E & $x \le -3$ oder $x \ge 1$ (ausserhalb; Grenzen eingeschlossen)\\
\end{raster}
\fehler{(a) Nur ein Fall ($x = 5$): Fälle- und Lösungs-Zeile 0 $\to$ 2 von 4\,P.
(b) $-3 \le x \le 1$ (zwischen statt ausserhalb): Lösungs-Zeile (b) 0.
(b) $x < -3$ oder $x > 1$ (Grenzen ausgeschlossen): Lösungs-Zeile (b) 0.}
\end{lpaufgabe}

\begin{lpaufgabe}{G7}{Lösungen zählen}{3 P}
\begin{raster}
$c = 2$ & 1\E & vier Lösungen: Die Waagrechte liegt unter dem Buckel $(1 \mid 4)$ und schneidet alle vier Äste des W.\\
$c = 4$ & 1\E & drei Lösungen: Sie berührt den Buckel ($x = 1$) und schneidet die äusseren Äste.\\
$c = 5$ & 1\E & zwei Lösungen: nur noch die äusseren Äste.\\
\end{raster}
\fehler{Gleichwertig zur Skizze ist die Rechnung $x^2 - 2x - 3 = \pm c$, also $(x - 1)^2 = 4 \pm c$: zwei, eine oder keine Lösung je Fall.
Nur die drei Zahlen, ohne Skizze und ohne Rechnung: alle Zeilen 0 — (E) verlangt einen erkennbaren Weg.
Folgefehler aus G4 (falscher Buckel): richtig gezählt zur eigenen Skizze gibt die Punkte nicht — (E).}
\end{lpaufgabe}

\begin{lpaufgabe}{G5}{Abschnittsweise schreiben}{4 P}
\begin{raster}
(a) Grenze & 1 & $2x - 5 = 0 \Rightarrow x = 2.5$\\
(a) Terme & 1\E & $|2x - 5| = 2x - 5$ für $x \ge 2.5$, $-2x + 5$ für $x < 2.5$. Beide nötig.\\
(b) Boden & 1\E & Höhe $4 - (-1) = 5$, von $x = -1$ bis $x = 4$. Höhe und Bereich nötig.\\
(b) Wert & 1\E & $y(6) = |7| + |2| = 9$\\
\end{raster}
\fehler{(a) Links $2x + 5$ (nur eine Zahl umgedreht): Terme-Zeile 0.
(a) Grenze bei $5$: Grenz-Zeile 0; die Terme mit falscher Grenze ebenfalls 0.
(b) Höhe $3$ (aus $|{-1} + 4|$ statt aus dem Abstand der Stellen): Boden-Zeile 0.}
\end{lpaufgabe}

\begin{lpaufgabe}{G6}{Gleichung und Ungleichung}{4 P}
\begin{raster}
(a) Fälle & 1 & $3x - 6 = 9$ oder $3x - 6 = -9$ (gleichwertig: $|x - 2| = 3$)\\
(a) Lösung & 1\E & $L = \{-1;\, 5\}$\\
(b) Schnittstellen & 1 & $x - 2 = \pm 3$: $x = -1$, $x = 5$\\
(b) Lösung & 1\E & $-1 < x < 5$ (zwischen den Schnittstellen, Grenzen ausgeschlossen)\\
\end{raster}
\fehler{(a) Nur ein Fall ($x = 5$): Fälle- und Lösungs-Zeile 0 $\to$ 2 von 4\,P.
(b) $x < 5$ oder «$x < -1$ oder $x > 5$»: Lösungs-Zeile (b) 0.
(b) $-1 \le x \le 5$ (Grenzen eingeschlossen): Lösungs-Zeile (b) 0.}
\end{lpaufgabe}

\begin{lpaufgabe}{G7}{Lösungen zählen}{3 P}
\begin{raster}
$c = 0.5$ & 1\E & vier Lösungen: Die Waagrechte liegt unter dem Buckel $(0 \mid 1)$ und schneidet alle vier Äste des W.\\
$c = 1$ & 1\E & drei Lösungen: Sie berührt den Buckel ($x = 0$) und schneidet die äusseren Äste ($x = \pm\sqrt2$).\\
$c = 3$ & 1\E & zwei Lösungen: nur noch die äusseren Äste ($x = \pm 2$).\\
\end{raster}
\fehler{Nur die drei Zahlen, ohne Skizze und ohne Begründung: alle Zeilen 0 — (E) verlangt einen erkennbaren Weg.
Buckel in der falschen Höhe (z.\,B. bei $-1$ belassen): die betroffenen Zeilen 0.}
\end{lpaufgabe}

\needspace{14\baselineskip}
\section*{4 · Selbsteinschätzung}
\begin{tabularx}{\linewidth}{@{}l X@{}}\toprule
22 – 25 P & Die geprüften Teile sitzen. Wo du Punkte verloren hast: das Kapitel dieser Aufgabe nochmals (Zuordnung unten).\\
17 – 21 P & Den schwächsten Teil nochmals: Simulation und Übungen des Kapitels, in dem du die meisten Punkte verloren hast.\\
11 – 16 P & Zurück zu den Kapiteln aller Aufgaben, in denen du Punkte verloren hast.\\
0 – 10 P & Zurück zu Kapitel 1 und von dort der Reihe nach weiter.\\\midrule
\multicolumn{2}{@{}p{\linewidth}@{}}{\small\textbf{Aufgabe $\to$ Kapitel:} G1 $\to$ Kapitel 1 · G2, G3 $\to$ Kapitel 2 · G4 $\to$ Kapitel 3 · G5 $\to$ Kapitel 4 · G6, G7 $\to$ Kapitel 5}\\\bottomrule
\end{tabularx}
\end{document}
